\DocumentMetadata{
  pdfversion=2.0,
  pdfstandard=ua-2,
  lang=en-US,
  tagging=on,
  tagging-setup={math/setup=mathml-SE,tabsorder=structure}
}
\documentclass[11pt,letterpaper]{article}
\usepackage[margin=0.68in,includefoot]{geometry}
\usepackage{newcomputermodern}
\usepackage{amsmath,amscd,mathtools}
\usepackage{tikz,adjustbox}
\providecommand{\square}{\mdlgwhtsquare}
\usepackage[hidelinks]{hyperref}
\hypersetup{
  pdftitle={Algebra First-Year Exam, May 2018, Part 1},
  pdfauthor={University of Florida Department of Mathematics},
  pdfsubject={Graduate mathematics examination},
  pdfdisplaydoctitle=true,
  bookmarksopen=true
}
\setcounter{secnumdepth}{0}
\setlength{\parindent}{0pt}
\setlength{\parskip}{0.28em}
\setlength{\emergencystretch}{3em}
\setlength{\leftmargini}{2.15em}
\setlength{\leftmarginii}{2.65em}
\setlength{\leftmarginiii}{3em}
\pagestyle{empty}
\raggedbottom
\allowdisplaybreaks
\NewTaggingSocketPlug{block/list/label}{examproblem}{%
  \tagstructbegin{tag=Lbl}\tagstructend%
  \tagstructbegin{tag=\LBody}%
  \tagstructbegin{tag=\ProblemHeadingTag,title-o={\ProblemHeadingText}}%
  \tagmcbegin{tag=\ProblemHeadingTag,actualtext={\ProblemHeadingText}}%
  #2%
  \tagmcend%
  \tagstructend%
}
\newenvironment{examproblems}[2]{%
  \def\ProblemHeadingTag{#2}%
  \AssignTaggingSocketPlug{block/list/label}{examproblem}%
  \begin{enumerate}%
  \setcounter{enumi}{#1}%
  \renewcommand{\labelenumi}{\textbf{\theenumi.}}%
  \setlength{\topsep}{0.35em}%
  \setlength{\partopsep}{0pt}%
  \setlength{\itemsep}{0.48em}%
  \setlength{\parsep}{0.18em}%
}{\end{enumerate}}
\newenvironment{examsubparts}{%
  \AssignTaggingSocketPlug{block/list/label}{default}%
  \begin{enumerate}%
  \setlength{\topsep}{0.18em}%
  \setlength{\partopsep}{0pt}%
  \setlength{\itemsep}{0.16em}%
  \setlength{\parsep}{0.1em}%
}{\end{enumerate}}
\newenvironment{examinstructions}{%
  \par\begingroup\itshape
}{%
  \par\endgroup\vspace{0.18em}
}
\makeatletter
\renewcommand\subsection{\@startsection{subsection}{2}{0pt}%
  {-1.25ex plus -.25ex minus -.1ex}%
  {0.55ex plus .1ex}%
  {\normalfont\large\bfseries}}
\renewcommand{\@maketitle}{%
  \newpage\null\vskip -1em%
  {\footnotesize\bfseries
    \href{https://www.ufl.edu/}{University of Florida}, \href{https://math.ufl.edu/}{Department of Mathematics}\par}%
  \vskip -0.5em%
  \tagstructbegin{tag=H1,title={Algebra First-Year Exam, May 2018, Part 1}}%
  \tagpdfparaOff%
  \tagmcbegin{tag=H1}%
    {\LARGE\bfseries\href{https://gma.math.ufl.edu/exams/algebra-first-year/}{Algebra First-Year Exam}, May 2018, Part 1\par}%
  \tagmcend%
  \tagpdfparaOn%
  \tagstructend%
  \vskip 0.25em%
  \tagmcbegin{artifact}%
    \hrule height 1pt%
  \tagmcend%
  \vskip 1em%
}
\makeatother
\title{Algebra First-Year Exam, May 2018, Part 1}
\author{}
\date{}
\begin{document}
\maketitle
\thispagestyle{empty}
\begin{examinstructions}
Answer four problems. You should indicate which problems you wish to have graded. Write your answers clearly in complete English sentences. You may quote results (within reason) as long as you state them clearly.
\end{examinstructions}
\begin{examproblems}{0}{H2}
\def\ProblemHeadingText{Problem 1.}
\item
Let~\(G\) be a finite group and let \(x,y \in G\) satisfy \(|x|=m\) and \(|y|=n\), with \(\gcd(m,n)=1\). (Here \(|x|\) denotes the order of~\(x\).)
\begin{examsubparts}
\item[\textnormal{(a)}]
Prove that if \(xy=yx\) and \(\gcd(m,n)=1\) then \(|xy|=mn\).
\item[\textnormal{(b)}]
Give an example where \(xy \ne yx\), \(\gcd(m,n)=1\), and \(|xy|>mn\).
\item[\textnormal{(c)}]
Give an example where \(xy \ne yx\), \(\gcd(m,n)=1\), and \(|xy|<mn\).
\end{examsubparts}
\def\ProblemHeadingText{Problem 2.}
\item
Let~\(G\) be a group and let \(S \subset G\). Let~\(N\) be the intersection of all normal subgroups~\(H\) of~\(G\) such that \(S \subset H\).
\begin{examsubparts}
\item[\textnormal{(a)}]
Prove that \(N \trianglelefteq G\).
\item[\textnormal{(b)}]
Prove that \(N=\langle T\rangle\), where \(T=\{gxg^{-1}:g\in G,\ x\in S\}\).
\end{examsubparts}
\def\ProblemHeadingText{Problem 3.}
\item
This problem has two parts.
\begin{examsubparts}
\item[\textnormal{(a)}]
Let \(n \ge 3\) and let \(\sigma,\tau \in S_n\). Prove that~\(\sigma\) and~\(\tau\) are conjugate in \(S_n\) if and only if~\(\sigma\) and~\(\tau\) have the same cycle type.
\item[\textnormal{(b)}]
Give an example of \(\sigma,\tau \in A_5\) such that~\(\sigma\) and~\(\tau\) are conjugate in \(S_5\) but~\(\sigma\) and~\(\tau\) are not conjugate in \(A_5\). Prove that your example is correct.
\end{examsubparts}
\def\ProblemHeadingText{Problem 4.}
\item
Let~\(G\) be a group of order \(715=5\cdot 11\cdot 13\). Prove that \(Z(G)\) is nontrivial.
\def\ProblemHeadingText{Problem 5.}
\item
Determine the number of abelian groups of order \(324=2^2\cdot 3^4\). Give a representative for each isomorphism class of such groups.
\def\ProblemHeadingText{Problem 6.}
\item
This problem has two parts.
\begin{examsubparts}
\item[\textnormal{(a)}]
Let~\(H\) and~\(K\) be groups and let \(\phi:K\to \operatorname{Aut}(H)\) be a homomorphism. Give the definition of the semidirect product \(H\rtimes_\phi K\). You should give the group operation on \(H\rtimes_\phi K\), but you don’t need to prove that the group axioms are satisfied.
\item[\textnormal{(b)}]
Give an example of abelian groups~\(H\) and~\(K\), a homomorphism \(\phi:K\to \operatorname{Aut}(H)\), and two elements of \(H\rtimes_\phi K\) which don’t commute.
\end{examsubparts}
\end{examproblems}
\end{document}
